\frontmatter%前言
\section*{前言}
\addcontentsline{toc}{chapter}{前言}

信号与系统对电子信息等相关专业的重要性是不言而喻的。本书基于以下四个目的进行编写，可供本科二年级以上的学生以及相关课程的授课教师参考。

第一，现有教材均在数学上做了简化，而本书结合以傅里叶分析为代表的数学专业课程，补全了这些内容，以供希望深入学习信号与系统，或者数学专业的读者参考。

第二，本书补充了大量课程以外的内容，它们或为知识的数学背景，或为课内知识在类比推理下的自然延伸，或为后续课程内容的简要介绍。这些内容旨在帮助读者节省查找资料的时间，快速建立对此类内容涉及到的问题及其背景的直观认识。

第三，许多工科教材的行文逻辑都不是完整的，这或许是出于读者普遍对所讨论的问题背景有所了解。本书力求给每个问题、定义和定理给出严谨、完整的说明，还将探索定理的边界条件，减小初学者和其他专业读者的阅读障碍。

第四，国内外的信号与系统课程以及后续相关课程对于傅里叶变换的定义有所不同，一种为频率形式，另一种为角频率形式，这种定义上的不同会导致诸多公式的内核相同但形式不同。本书将在介绍傅里叶变换时同时给出两种形式的公式。我们将看到，同时给出两种傅里叶变换，不仅方便查找、对比公式，理解傅里叶变换的本质，还能在必要时转向其中一种傅里叶变换，使得讨论更加简洁明快。

为了使初学者也能够参考本书，一些小节和定理标注了“*”以提示读者可以选择性地阅读。

本书第一章的“线性代数”一节由北京市第四中学2025届毕业生王怀宇完成，他也为本书的其他部分提出了许多有用的意见。本书的其余部分由严子竣完成。

限于作者水平，书中难免存在不足之处，希望读者可以提出宝贵的意见或建议，可以反馈至我的电子邮箱：2024212622@bupt.cn。

{
\raggedleft{}
北京市十一学校2024届毕业生\\
严子竣\\
2026年春\\
}

\section*{致谢}
\addcontentsline{toc}{chapter}{致谢}

感谢北京邮电大学电子工程学院尹霄丽教授对本书编写工作的指导与支持。感谢所有阅读本书初稿并提出修改意见的同学。
